\specsection{The theory of Brownian motion}

This theory occupies a special position in contemporary scientific development, by the constant and fertile exchange between physical problems and `pure' mathematics to which it bears witness. The study of Brownian motion, discovered in 1829 by the botanist Brown, had been conducted intensively in the 19th century by numerous physicists, \(^{(3)}\) but the first satisfactory mathematical model was only invented by Einstein in 1905. In the simple case of a particle moving along a straight line, Einstein's fundamental hypotheses may be formulated as follows: if \(x(t)\) is the abscissa of the particle at the instant t, and if \(t_{0} < t_{1} < \cdots < t_{n-1} < t_{n}\) , the successive displacements \(x(t_{i}) - x(t_{i-1})\) (for \(1 \leqslant i \leqslant n\) ) are independent Gaussian random variables. This is not the place to evoke in detail the important experimental work of J. Perrin that motivated Einstein's theory; for our purposes, we need only retain a remark of Perrin, according to which the observation of trajectories of Brownian motion irresistibly suggested to him the ``mathematicians' functions without derivative''. This remark was to be the initial spark for Wiener.

Quite another current of ideas has its origins in the kinetic theory of gases, developed between 1870 and 1900 by Boltzmann and Gibbs. Consider a gas formed by N molecules of mass m at (absolute) temperature T, and denote by \(v_{1},\ldots,v_{N}\) the velocities of the N molecules of the gas; the kinetic energy of the system is equal to

\[
\frac {m}{2} (\mathbf {v} _ {1} ^ {2} + \dots + \mathbf {v} _ {\mathrm{N}} ^ {2}) = 3 \mathrm{N} k \mathrm{T},\tag{1}
\]

where k is Boltzmann's constant. According to the ideas of Gibbs, the multitude of shocks between molecules does not allow the precise determination of the velocities of the molecules, and it is convenient to introduce a probability law P on the sphere S of the space of dimension 3N defined by the equation (1). The `micro-canonical' hypothesis consists in assuming that P is the measure on the sphere S with mass 1 invariant under rotation. Moreover, Maxwell's law of velocities states that the law of probability of a component of the velocity of a molecule is a Gaussian measure with variance \(2k\mathrm{T} / m\) (§6, No.~5, Remark 3). Borel seems to have been the first to observe in 1914 that Maxwell's law is a consequence of Gibbs' hypotheses and the properties of the sphere when the number of molecules is very large. He considered a sphere S in a Euclidean space of large dimension and the measure P of mass 1 on S invariant under rotation; using the classical approximation methods based on Stirling's formula, he showed that the projection of P on a coordinate axis is approximately Gaussian. These results were sharpened a little later by Gâteaux and Lévy (IX, \(a\) ). Given an integer \(m \geqslant 1\) and a number \(r > 0\) , denote by \(\mathrm{S}_{m,r}\) the set of sequences of the form \((x_{1}, \ldots, x_{m}, 0, 0, \ldots)\) with \(x_{1}^{2} + \cdots + x_{m}^{2} = r^{2}\) ; also, denote by \(\sigma_{m,r}\) the measure with mass 1 on \(\mathrm{S}_{m,r}\) invariant under rotation. Stated in modern language, the result of Gâteaux and Lévy is as follows: the sequence of measures \(\sigma_{m,1}\) tends tightly to a unit mass at the origin \((0, 0, \ldots)\) , and the sequence of measures \(\sigma_{m,\sqrt{m}}\) tends tightly to a measure \(\Gamma\) of the form

\[
d \Gamma (x _ {1}, x _ {2}, \dots) = \prod_ {n = 1} ^ {\infty} d \gamma (x _ {n})
\]

(γ is the Gaussian measure on R with variance 1).

The preceding measure \(\Gamma\) plays the role of a Gaussian measure in infinite dimensions. It seems indeed that Lévy had confusedly hoped to define in an intrinsic manner a Gaussian measure on every infinite-dimensional Hilbert space. In fact, as shown by Lévy and Wiener, the measure \(\Gamma\) is invariant in a certain sense \(^{(4)}\) under the automorphisms of \(l^{2}\) ; unfortunately, the set \(l^{2}\) of square-summable sequences \((x_{1}, x_{2}, \ldots, x_{n}, \ldots)\) has measure zero for \(\Gamma\) . It is now known that one must be content to have a Gaussian promeasure on an infinite-dimensional Hilbert space. \(^{(5)}\)

We owe to Wiener the essential progress: if one does not have a reasonable Gaussian measure on an infinite-dimensional Hilbert space, one can construct by the operation of primitive a measure \(w\) on a space of continuous functions starting from a Gaussian promeasure (cf.~§6, No.~7, Th. 1 for the details). We shall explain succinctly Wiener's original construction of \(w\) (X); it is directly influenced by the relation \(\Gamma = \lim_{m \to \infty} \sigma_{m,\sqrt{m}}\) of Gâteaux and Lévy. For every integer \(m \geqslant 1\) , denote by \(H_m\) the set of functions on \(T = ]0,1]\) that are constant on each of the intervals \(\left[ \frac{k-1}{m}, \frac{k}{m} \right]\) (for \(k = 1,2,\ldots,m\) ), and by \(\pi_m\) the measure of mass 1, invariant under rotation, on the Euclidean sphere of radius 1 in \(R^m\) . Let \(f_m\) be the isomorphism of \(H_m\) onto \(R^m\) that associates, to each function taking the value \(a_k\) on the interval \(\left[ \frac{k-1}{m}, \frac{k}{m} \right]\) , the vector \((a_1, a_2 - a_1, \ldots, a_m - a_{m-1})\) (whence the term `differential space' dear to Wiener); denote by \(w_m\) the measure on \(H_m\) that is the image of \(\pi_m\) under \(f_m^{-1}\) . Wiener defined the desired measure \(w\) as the limit of the measures \(w_m\) . To be precise, let us denote by \(H\) the set of regulated functions on \(T\) , with the topology of uniform convergence (we have \(H_m \subset H\) for every integer \(m \geqslant 1\) ); for every bounded uniformly continuous function \(F\) on \(H\) , the limit \(A\{F\} = \lim_{m \to \infty} \int_{H_m} F(x) dw_m(x)\) exists; next, Wiener obtained certain upper bounds by a subtle analysis of the fluctuations of the game of heads-or-tails, and taking up again the compactness arguments highlighted by Daniell, he showed that one is under the conditions for applying Daniell's extension theorem. One concludes the existence of a measure \(w\) carried by \(C(T)\) and such that \(A\{F\} = \int_{C(T)} F(x) dw(x)\) . Wiener was then able to show that the measure \(w\) corresponds to Einstein's hypotheses, \(^{(6)}\) and his estimates allowed him to give a precise meaning to Perrin's remark on functions without derivatives: the set of functions satisfying a Lipschitz condition of order \(\frac{1}{2}\) is negligible for \(w\) (however, for every \(a\) with \(0 < a < \frac{1}{2}\) , almost every function satisfies a Lipschitz condition of order \(a\) ).

Today, numerous constructions of the Wiener measure are known. Thus, Paley and Wiener use random Fourier series (XI, Ch. IX): for every real

\[
\begin{array}{l} \int_ {\mathcal {C} (\mathrm{T})} f (x (t _ {1}), \dots , x (t _ {n})) d w (x) = \\ (2 \pi) ^ {- n / 2} \prod_ {i = 1} ^ {n} (t _ {i} - t _ {i - 1}) ^ {- 1 / 2} \int \dots \int f (x _ {1}, \dots , x _ {n}) \exp \left(- \frac {1}{2} \sum_ {i = 1} ^ {n} \frac {(x _ {i} - x _ {i - 1}) ^ {2}}{t _ {i} - t _ {i - 1}}\right) d x _ {1} \dots d x _ {n}, \end{array}
\]

sequence \(\mathbf{a} = (a_n)_{n\geqslant 1}\) and every integer \(m\geqslant 0\) , let us define the function \(f_{m,\mathbf{a}}\) on ]0,1] by

\[
f _ {m, \mathbf {a}} (t) = a _ {1} t + 2 \sum_ {k = 2} ^ {2 ^ {m + 1}} \frac {1}{\pi k} a _ {k - 1} \sin \pi k t;
\]

one can show that for \(\Gamma\) -almost every sequence \(\mathbf{a}\) , the sequence of functions \(f_{m,\mathbf{a}}\) tends to a continuous function \(f_{\mathbf{a}}\) , and that \(w\) is the image of \(\Gamma\) under the mapping (defined almost everywhere) \(\mathbf{a} \mapsto f_{\mathbf{a}}\) . Later on, Lévy gave in (IX, \(b\) ), \(c\) ) a construction very near to that which we have exposed in §6, No.~7. Finally, Kac, Donsker and Erdős showed around 1950 how to replace the spherical measures \(\pi_m\) on \(\mathbf{R}^m\) in Wiener's original construction by more general measures. Their results establish a solid link between Wiener measure and the limit theorems of Probability Theory; they were to be completed and systematized by Prokhorov (XIII) in a work to which we shall return later on.

This is not the place to analyze the numerous and important probabilistic works brought about by Wiener's discovery; nowadays, Brownian motion appears only as one of the most important examples of a Markoff process. We shall only mention Kac's application of Wiener measure to the solution of certain parabolic partial differential equations; this is a question of adapting ideas of Feynmann in quantum field theory---yet another example of the reciprocal influence of mathematics and problems of physics.

